\subsection{bin-fold escort 族的 Blackwell 等价、圆弧几何与决策塌缩}\label{subsec:conclusion-binfold-escort-blackwell-geometry-risk}
在定理 \ref{thm:conclusion-binfold-escort-full-fdiv-collapse} 已把 escort--escort 的全部 \(f\)-散度压缩为 Bernoulli 两点闭式、定理 \ref{thm:conclusion-escort-fisher-rao-geodesic-reparametrization} 已给出 Fisher--Rao 距离的完全显式表示、而定理 \ref{thm:conclusion-chi2-linear-recovers-phi} 已说明单个 \(\chi^2\)-型极限常数能够线性恢复黄金参数之后，还可把这条主线进一步提升到统计实验与决策理论层面。其核心事实是：末位读出并非仅提供一个渐近近似，而是以指数精度给出整个 escort 温度族的显式 Blackwell 代表；因此，全 \(f\)-散度、其总变差/\(\chi^2\) 截面、Fisher 曲率以及有界决策风险都塌缩到同一 Bernoulli 两点骨架。

\begin{theorem}[末位读出核给出的显式渐近 Blackwell 等价]\label{thm:conclusion-binfold-escort-blackwell-equivalence}
\leanverified{paper\_conclusion\_binfold\_escort\_blackwell\_equivalence}
固定紧集 \(\Theta\subset[0,\infty)\)。对每个 \(m\ge 1\)，定义末位两扇区
\[
X_{m,0}:=\{w\in X_m:\ w_m=0\},
\qquad
X_{m,1}:=\{w\in X_m:\ w_m=1\},
\]
读出核
\[
K_m:X_m\to\{0,1\},
\qquad
K_m(w):=w_m,
\]
以及回灌核
\[
R_m:\{0,1\}\rightsquigarrow X_m,
\qquad
R_m(0):=\mathrm{Unif}(X_{m,0}),
\qquad
R_m(1):=\mathrm{Unif}(X_{m,1}).
\]
再记
\[
\theta(q):=\frac{1}{1+\varphi^{q+1}},
\qquad
\beta_q:=\mathrm{Bernoulli}(\theta(q)).
\]
则有一致总变差估计
\[
\sup_{q\in\Theta}
\bigl\|
\pi_{m,q}-R_m\sharp\beta_q
\bigr\|_{\mathrm{TV}}
=
O_\Theta\!\left(\Bigl(\frac{\varphi}{2}\Bigr)^m\right),
\]
以及
\[
\sup_{q\in\Theta}
\bigl\|
K_m\sharp\pi_{m,q}-\beta_q
\bigr\|_{\mathrm{TV}}
=
O_\Theta\!\left(\Bigl(\frac{\varphi}{2}\Bigr)^m\right).
\]
因此，对实验族
\[
\mathcal E_m(\Theta):=\{\pi_{m,q}:q\in\Theta\},
\qquad
\mathcal B(\Theta):=\{\beta_q:q\in\Theta\},
\]
有
\[
\Delta_{\mathrm{LeCam}}\bigl(\mathcal E_m(\Theta),\mathcal B(\Theta)\bigr)
=
O_\Theta\!\left(\Bigl(\frac{\varphi}{2}\Bigr)^m\right).
\]
特别地，末位统计量 \(w\mapsto w_m\) 构成 bin-fold escort 温度族的渐近充分统计量。
\end{theorem}
\begin{proof}
取 \(Q\ge 0\) 使 \(\Theta\subset[0,Q]\)。由定理 \ref{thm:killo-fold-bin-escort-onebit-fisher}，对所有充分大的 \(m\) 有
\[
\sup_{q\in\Theta}\|\pi_{m,q}-\nu_{m,q}\|_{\mathrm{TV}}
=
O_\Theta\!\left(\Bigl(\frac{\varphi}{2}\Bigr)^m\right),
\]
其中
\[
\nu_{m,q}(w)
:=
\frac{u_m(w)\varphi^{-q w_m}}{Z_{m,q}},
\qquad
Z_{m,q}:=\frac{F_{m+1}+\varphi^{-q}F_m}{F_{m+2}}.
\]
把
\[
\theta_m(q):=\nu_{m,q}(X_{m,1})
=
\frac{\varphi^{-q}F_m}{F_{m+1}+\varphi^{-q}F_m}
\]
记为有限尺度末位质量，则由 \(\nu_{m,q}\) 的块内均匀性可知
\[
\nu_{m,q}=R_m\sharp\mathrm{Bernoulli}(\theta_m(q)).
\]
另一方面，
\[
\theta_m(q)=\frac{\varphi^{-q}}{F_{m+1}/F_m+\varphi^{-q}},
\qquad
\theta(q)=\frac{\varphi^{-q}}{\varphi+\varphi^{-q}}.
\]
由于 \(F_{m+1}/F_m=\varphi+O(\varphi^{-2m})\)，可得
\[
\sup_{q\in\Theta}|\theta_m(q)-\theta(q)|
=
O_\Theta(\varphi^{-2m})
=
O_\Theta\!\left(\Bigl(\frac{\varphi}{2}\Bigr)^m\right).
\]
于是
\[
\sup_{q\in\Theta}
\|\nu_{m,q}-R_m\sharp\beta_q\|_{\mathrm{TV}}
\le
\sup_{q\in\Theta}|\theta_m(q)-\theta(q)|
=
O_\Theta\!\left(\Bigl(\frac{\varphi}{2}\Bigr)^m\right),
\]
再与 \(\|\pi_{m,q}-\nu_{m,q}\|_{\mathrm{TV}}\) 的一致估计结合，即得第一式。

对第二式，先注意到
\[
K_m\sharp\nu_{m,q}
=
\mathrm{Bernoulli}(\theta_m(q)).
\]
再由推前下总变差不增，
\[
\|K_m\sharp\pi_{m,q}-K_m\sharp\nu_{m,q}\|_{\mathrm{TV}}
\le
\|\pi_{m,q}-\nu_{m,q}\|_{\mathrm{TV}},
\]
并结合 \(|\theta_m(q)-\theta(q)|\) 的一致估计，即得
\[
\sup_{q\in\Theta}
\|K_m\sharp\pi_{m,q}-\beta_q\|_{\mathrm{TV}}
=
O_\Theta\!\left(\Bigl(\frac{\varphi}{2}\Bigr)^m\right).
\]

最后，两条显式核 \(K_m\) 与 \(R_m\) 分别给出双向 deficiency 的同阶上界，故 Le Cam 距离满足同阶估计。证毕。
\end{proof}

\begin{corollary}[平方根嵌入下的精确圆弧参数化]\label{cor:conclusion-binfold-escort-sqrt-circle-arc}
\leanverified{paper\_conclusion\_binfold\_escort\_sqrt\_circle\_arc}
定义平方根嵌入
\[
\Xi(q):=
\bigl(\sqrt{1-\theta(q)},\,\sqrt{\theta(q)}\bigr)\in S^1_+,
\qquad
\theta(q)=\frac{1}{1+\varphi^{q+1}}.
\]
再定义角变量
\[
\alpha(q):=\arctan\!\bigl(\varphi^{-(q+1)/2}\bigr)\in(0,\pi/2).
\]
则有精确表示
\[
\Xi(q)=\bigl(\cos\alpha(q),\,\sin\alpha(q)\bigr).
\]
因此 bin-fold escort 极限族在 Fisher--Rao 度量下恰是一段圆弧测地线，并且
\[
d_{\mathrm{FR}}(q,p)
=
2\bigl|\alpha(q)-\alpha(p)\bigr|
=
2\left|
\arctan\!\bigl(\varphi^{-(q+1)/2}\bigr)
-
\arctan\!\bigl(\varphi^{-(p+1)/2}\bigr)
\right|.
\]
与此同时，其 Fisher 信息密度满足
\[
\mathcal I(q)
=
(\log\varphi)^2\,\theta(q)\bigl(1-\theta(q)\bigr)
=
(\log\varphi)^2\frac{\varphi^{q+1}}{(1+\varphi^{q+1})^2}.
\]
\end{corollary}
\begin{proof}
信息密度与 Fisher--Rao 距离闭式已由定理 \ref{thm:conclusion-escort-fisher-rao-geodesic-reparametrization} 给出。只需说明圆弧参数化。
令
\[
u:=\varphi^{-(q+1)/2},
\]
则
\[
\theta(q)=\frac{u^2}{1+u^2},
\qquad
1-\theta(q)=\frac{1}{1+u^2}.
\]
因此
\[
\Xi(q)=\frac{(1,u)}{\sqrt{1+u^2}}
=
\bigl(\cos(\arctan u),\,\sin(\arctan u)\bigr),
\]
即
\[
\Xi(q)=\bigl(\cos\alpha(q),\,\sin\alpha(q)\bigr).
\]
再由定理 \ref{thm:conclusion-escort-fisher-rao-geodesic-reparametrization} 中
\[
d_{\mathrm{FR}}(q,p)
=
2\left|
\arcsin\sqrt{\theta(q)}-\arcsin\sqrt{\theta(p)}
\right|
\]
与恒等式
\[
\arcsin\sqrt{\theta(q)}=\arctan\!\bigl(\varphi^{-(q+1)/2}\bigr)
\]
即可得到所示另一种距离表示。证毕。
\end{proof}

\begin{corollary}[escort 温度射线的 Fisher 有限直径]\label{cor:conclusion-binfold-escort-fisher-finite-diameter}
\leanverified{paper\_conclusion\_binfold\_escort\_fisher\_finite\_diameter}
沿用推论 \ref{cor:conclusion-binfold-escort-sqrt-circle-arc} 的记号。则无界温度射线 \([0,\infty)\) 在 Fisher--Rao 度量下具有有限总长度
\[
\operatorname{Len}_{\mathrm{FR}}([0,\infty))
=
\sup_{q\ge 0}d_{\mathrm{FR}}(0,q)
=
2\arctan\!\bigl(\varphi^{-1/2}\bigr)
<\infty.
\]
更一般地，对任意 \(0\le p\le q<\infty\) 都有
\[
d_{\mathrm{FR}}(p,q)
=
2\left(
\arctan\!\bigl(\varphi^{-(p+1)/2}\bigr)
-\arctan\!\bigl(\varphi^{-(q+1)/2}\bigr)
\right).
\]
\end{corollary}
\begin{proof}
由推论 \ref{cor:conclusion-binfold-escort-sqrt-circle-arc}，
\[
d_{\mathrm{FR}}(p,q)
=
2\bigl|\alpha(p)-\alpha(q)\bigr|,
\qquad
\alpha(s)=\arctan\!\bigl(\varphi^{-(s+1)/2}\bigr).
\]
当 \(0\le p\le q\) 时，\(\alpha\) 严格递减，故绝对值可去掉，得到第二式。再令 \(p=0\) 并令 \(q\to\infty\)，由 \(\alpha(q)\to 0\) 可知
\[
\sup_{q\ge 0}d_{\mathrm{FR}}(0,q)=2\alpha(0)=2\arctan\!\bigl(\varphi^{-1/2}\bigr).
\]
这正是 \([0,\infty)\) 的 Fisher--Rao 总长度。证毕。
\end{proof}

\begin{proposition}[Csiszár 二态闭包、总变差/\texorpdfstring{$\chi^2$}{chi2} 截面与黄金识别]\label{prop:conclusion-binfold-escort-csiszar-blackwell-phi}
\leanverified{paper\_conclusion\_binfold\_escort\_csiszar\_blackwell\_phi}
同一两态结构同时控制 escort 温度族内部的全部 Csiszár--Morimoto 几何、其总变差/\(\chi^2\) 截面，以及相对均匀基线下的黄金参数识别。更具体地：
\begin{enumerate}
  \item 对任意连续凸函数 \(f:(0,\infty)\to\mathbb R\) 且 \(f(1)=0\)，以及任意 \(p,q\ge 0\)，定理 \ref{thm:conclusion-binfold-escort-full-fdiv-collapse} 给出
  \[
  \lim_{m\to\infty}D_f(\pi_{m,q}\Vert \pi_{m,p})
  =
  \theta(p)\,f\!\Bigl(\frac{\theta(q)}{\theta(p)}\Bigr)
  +
  \bigl(1-\theta(p)\bigr)\,
  f\!\Bigl(\frac{1-\theta(q)}{1-\theta(p)}\Bigr).
  \]
  \item 总变差极限满足
  \[
  \lim_{m\to\infty}\|\pi_{m,q}-\pi_{m,p}\|_{\mathrm{TV}}
  =
  |\theta(q)-\theta(p)|.
  \]
  \item 若记
  \[
  f_2(u):=(u-1)^2
  \]
  并定义
  \[
  \chi^2(\mu\Vert\nu):=D_{f_2}(\mu\Vert\nu),
  \]
  则 Pearson \(\chi^2\) 极限满足
  \[
  \lim_{m\to\infty}\chi^2(\pi_{m,q}\Vert \pi_{m,p})
  =
  \frac{(\theta(q)-\theta(p))^2}{\theta(p)(1-\theta(p))}.
  \]
  \item 对相对均匀基线的二原子极限，定理 \ref{thm:conclusion-chi2-linear-recovers-phi} 在同一
  \[
  f_2(u):=(u-1)^2
  \]
  截面上给出精确线性公式
  \[
  D_\infty^{f_2}=\frac{2\varphi-3}{5},
  \qquad
  \varphi=\frac{5D_\infty^{f_2}+3}{2}.
  \]
\end{enumerate}
因此，末位两态塌缩不仅把 escort--escort 的全部 \(f\)-散度压缩到同一条 Bernoulli 圆弧上，而且还把 TV/\(\chi^2\) 等经典截面封闭到同一单参数阴影，并允许由相对均匀基线上的单个 \(\chi^2\)-型极限标量对黄金参数作线性恢复。
\end{proposition}
\begin{proof}
第一条正是定理 \ref{thm:conclusion-binfold-escort-full-fdiv-collapse} 的结论。

第二条在第一条中取
\[
f(u)=\frac{|u-1|}{2}
\]
即可，因为该 \(f\)-散度恰为总变差距离，而 Bernoulli 两点模型上的总变差正等于参数差的绝对值。

第三条在第一条中取
\[
f_2(u)=(u-1)^2
\]
即可。此时
\[
\chi^2(\pi_{m,q}\Vert \pi_{m,p})
\longrightarrow
\bigl(1-\theta(p)\bigr)
\left(\frac{\theta(p)-\theta(q)}{1-\theta(p)}\right)^2
+\theta(p)
\left(\frac{\theta(q)-\theta(p)}{\theta(p)}\right)^2,
\]
化简便得
\[
\lim_{m\to\infty}\chi^2(\pi_{m,q}\Vert \pi_{m,p})
=
\frac{(\theta(q)-\theta(p))^2}{\theta(p)(1-\theta(p))}.
\]

第四条正是定理 \ref{thm:conclusion-chi2-linear-recovers-phi} 的结论。上述各式共同依赖于 bin-fold 退化谱的两态极限，因此可在同一识别链中统一表述。证毕。
\end{proof}

\begin{theorem}[有界决策问题的最优 Bayes 风险塌缩到 Bernoulli 极限]\label{thm:conclusion-binfold-escort-bayes-risk-collapse}
\leanverified{paper\_conclusion\_binfold\_escort\_bayes\_risk\_collapse}
固定紧参数集 \(\Theta\subset[0,\infty)\)，并令
\[
\mathcal E_m(\Theta):=\{\pi_{m,q}:q\in\Theta\},
\qquad
\mathcal B(\Theta):=\{\beta_q:q\in\Theta\},
\qquad
\beta_q=\mathrm{Bernoulli}(\theta(q)).
\]
设 \(\mathcal A\) 为任意标准 Borel 动作空间，\(\ell:\Theta\times\mathcal A\to\mathbb R\) 为有界损失函数，并记
\[
\|\ell\|_\infty:=\sup_{(q,a)\in\Theta\times\mathcal A}|\ell(q,a)|.
\]
对任意先验 \(\Pi\in\mathcal P(\Theta)\)，分别记
\[
\mathfrak R_m^\ast(\Pi,\ell)
\]
与
\[
\mathfrak R_\infty^\ast(\Pi,\ell)
\]
为实验 \(\mathcal E_m(\Theta)\) 与 \(\mathcal B(\Theta)\) 下的最优 Bayes 风险。则有
\[
\sup_{\Pi\in\mathcal P(\Theta)}
\left|
\mathfrak R_m^\ast(\Pi,\ell)-\mathfrak R_\infty^\ast(\Pi,\ell)
\right|
\le
2\|\ell\|_\infty\,
\Delta_{\mathrm{LeCam}}\bigl(\mathcal E_m(\Theta),\mathcal B(\Theta)\bigr),
\]
从而
\[
\sup_{\Pi\in\mathcal P(\Theta)}
\left|
\mathfrak R_m^\ast(\Pi,\ell)-\mathfrak R_\infty^\ast(\Pi,\ell)
\right|
=
O_\Theta\!\left(
\|\ell\|_\infty
\Bigl(\frac{\varphi}{2}\Bigr)^m
\right).
\]
换言之，对任意有界统计决策问题，bin-fold escort 实验族的极限最优 Bayes 风险完全由末位 Bernoulli 实验决定。
\end{theorem}
\begin{proof}
记
\[
\varepsilon_m:=
\Delta_{\mathrm{LeCam}}\bigl(\mathcal E_m(\Theta),\mathcal B(\Theta)\bigr).
\]
由定理 \ref{thm:conclusion-binfold-escort-blackwell-equivalence}，有
\[
\varepsilon_m=
O_\Theta\!\left(\Bigl(\frac{\varphi}{2}\Bigr)^m\right).
\]

先证
\(
\mathfrak R_m^\ast(\Pi,\ell)\le \mathfrak R_\infty^\ast(\Pi,\ell)+2\|\ell\|_\infty\varepsilon_m
\)。
取任意 \(\{0,1\}\) 上的决策核 \(\rho:\{0,1\}\rightsquigarrow\mathcal A\)，并构造 \(X_m\) 上的决策核
\[
\delta_m:=\rho\circ K_m.
\]
对任意 \(q\in\Theta\)，由推前下总变差不增与损失有界性，
\[
\left|
\int\ell(q,a)\,\delta_m\sharp\pi_{m,q}(\mathrm da)
-
\int\ell(q,a)\,\rho\sharp\beta_q(\mathrm da)
\right|
\le
2\|\ell\|_\infty\,
\|K_m\sharp\pi_{m,q}-\beta_q\|_{\mathrm{TV}}
\le
2\|\ell\|_\infty\varepsilon_m.
\]
再对 \(q\) 按先验 \(\Pi\) 积分，并对 \(\rho\) 取下确界，即得
\[
\mathfrak R_m^\ast(\Pi,\ell)\le \mathfrak R_\infty^\ast(\Pi,\ell)+2\|\ell\|_\infty\varepsilon_m.
\]

反向不等式同理。任取 \(X_m\) 上的决策核 \(\delta_m:X_m\rightsquigarrow\mathcal A\)，定义 \(\{0,1\}\) 上的决策核
\[
\rho:=\delta_m\circ R_m.
\]
则对任意 \(q\in\Theta\)，有
\[
\left|
\int\ell(q,a)\,\rho\sharp\beta_q(\mathrm da)
-
\int\ell(q,a)\,\delta_m\sharp\pi_{m,q}(\mathrm da)
\right|
\le
2\|\ell\|_\infty\,
\|R_m\sharp\beta_q-\pi_{m,q}\|_{\mathrm{TV}}
\le
2\|\ell\|_\infty\varepsilon_m.
\]
对先验积分并对 \(\delta_m\) 取下确界，得到
\[
\mathfrak R_\infty^\ast(\Pi,\ell)\le \mathfrak R_m^\ast(\Pi,\ell)+2\|\ell\|_\infty\varepsilon_m.
\]

合并两式即得
\[
\left|
\mathfrak R_m^\ast(\Pi,\ell)-\mathfrak R_\infty^\ast(\Pi,\ell)
\right|
\le
2\|\ell\|_\infty\varepsilon_m.
\]
再对 \(\Pi\in\mathcal P(\Theta)\) 取上确界，并代入 \(\varepsilon_m\) 的衰减率，结论成立。证毕。
\end{proof}

\begin{theorem}[单个幂散度极限标量对 escort 热力学阶梯的完全恢复]\label{thm:conclusion-powerdiv-scalar-recovers-escort-thermo-ladder}
\leanverified{paper\_conclusion\_powerdiv\_scalar\_recovers\_escort\_thermo\_ladder}
固定 \(\alpha>1\)，并记
\[
f_\alpha(t):=t^\alpha-1-\alpha(t-1),
\qquad
D_\alpha^\infty:=D_{f_\alpha}^\infty.
\]
则单个实数 \(D_\alpha^\infty\) 唯一决定黄金参数 \(\varphi\)。因而它也唯一决定整条二态参数曲线
\[
\theta(q)=\frac{1}{1+\varphi^{q+1}}
\qquad (q\ge 0),
\]
并进一步唯一决定 bin-fold escort 族的整个渐近热力学阶梯。更具体地，对任意 \(p,q\ge 0\)，都有
\[
\lim_{m\to\infty}D_{\mathrm{KL}}(\pi_{m,q}\Vert \pi_{m,p})
=
\theta(q)\log\frac{\theta(q)}{\theta(p)}
+
\bigl(1-\theta(q)\bigr)\log\frac{1-\theta(q)}{1-\theta(p)},
\]
以及
\[
Z_m^{(q)}\Rightarrow Z^{(q)},
\qquad
Z^{(q)}\in\{1,\varphi^{-1}\},
\]
\[
\mathbb P\bigl(Z^{(q)}=\varphi^{-1}\bigr)=\theta(q),
\qquad
\mathbb P\bigl(Z^{(q)}=1\bigr)=1-\theta(q).
\]
换言之，任意一个固定 \(\alpha>1\) 上的 \(f_\alpha\)-极限标量，已经足以恢复整个 escort 温度族的渐近相对熵几何与两尺度残差律。
\end{theorem}
\begin{proof}
命题 \ref{prop:fold-bin-phi-identifiable-power-divergence} 已给出
\[
D_\alpha^\infty
=
\frac{1}{5^{\alpha/2}}
\bigl(\varphi^{2\alpha-1}+\varphi^{\alpha-2}\bigr)-1
\]
并证明对每个固定 \(\alpha>1\)，上式右端关于 \(\varphi>1\) 严格单调，故 \(D_\alpha^\infty\) 唯一决定 \(\varphi\)。于是
\[
\theta(q)=\frac{1}{1+\varphi^{q+1}}
\]
对一切 \(q\ge 0\) 也随之唯一确定。

随后，定理 \ref{thm:conclusion-escort-two-state-closure} 立即给出所示的 KL 极限与两点残差极限。由于这些极限对象完全由 \(\theta(q)\) 与 \(\varphi^{-1}\) 描述，故它们都由单个标量 \(D_\alpha^\infty\) 唯一锁定。证毕。
\end{proof}

\begin{theorem}[单一幂散度的渐近统计完备性]\label{thm:conclusion-binfold-single-powerdiv-asymptotic-statistical-completeness}
\leanverified{paper\_conclusion\_binfold\_single\_powerdiv\_asymptotic\_statistical\_completeness}
在 bin-fold 主线中，固定任意 \(\alpha>1\)，并记
\[
f_\alpha(t):=t^\alpha-1-\alpha(t-1).
\]
设两组稳定层输出分布 \((\mu_m)\)、\((\mu'_m)\) 分别相对于其稳定层均匀基线 \(u_m,u'_m\) 满足
\[
\lim_{m\to\infty}D_{f_\alpha}(\mu_m\Vert u_m)
=
\lim_{m\to\infty}D_{f_\alpha}(\mu'_m\Vert u'_m).
\]
则二者具有同一黄金参数 \(\varphi\)，并因而同时具有：
\[
\text{同一全 Rényi 极限族 }\{D_\beta^\infty\}_{\beta>0},
\qquad
\text{同一全 Csisz\'ar 极限锥 }\{D_f^\infty\},
\]
\[
\text{同一两点残差律 }Z^{(q)}\in\{1,\varphi^{-1}\},
\qquad
\text{同一 escort Fisher 几何},
\qquad
\text{同一涨落--响应律}.
\]
更强地，这两组实验在渐近 Le Cam 意义下属于同一统计实验。
\end{theorem}

\begin{proof}
由定理 \ref{thm:conclusion-powerdiv-scalar-recovers-escort-thermo-ladder}，任意一个固定 \(\alpha>1\) 上的极限标量
\[
D_{f_\alpha}^\infty
\]
已足以唯一恢复黄金参数 \(\varphi\)。故一旦两组分布满足同一极限幂散度，它们便具有相同的 \(\varphi\)。

随后，命题 \ref{prop:conclusion-binfold-escort-csiszar-blackwell-phi}、定理 \ref{thm:conclusion-powerdiv-scalar-recovers-escort-thermo-ladder} 与定理 \ref{thm:conclusion-binfold-escort-bayes-risk-collapse} 共同说明：全 Rényi 极限族、全 Csisz\'ar 极限锥、两点残差律、escort Fisher 几何以及涨落--响应常数，全部都是 \(\varphi\) 的显式函数。故这些极限对象在两组实验间完全一致。

最后，由定理 \ref{thm:conclusion-binfold-escort-blackwell-equivalence}，两组 bin-fold 实验族都以指数精度塌缩到同一个由 \(\varphi\) 唯一决定的 Bernoulli 两点极限实验。因而它们分别与同一极限实验的 Le Cam 距离趋于零，进而彼此的 Le Cam 距离也趋于零。故它们属于同一渐近统计实验。证毕。
\end{proof}

\begin{corollary}[bin-fold 极限实验不存在非平凡 nuisance 变形]\label{cor:conclusion-binfold-no-nontrivial-nuisance-deformation-fixed-powerdiv}
在定理 \ref{thm:conclusion-binfold-single-powerdiv-asymptotic-statistical-completeness} 的设定下，任何保持某个固定 \(\alpha>1\) 的极限幂散度
\[
D_{f_\alpha}^\infty
\]
不变的变形，都是渐近统计上平凡的；亦即，它不可能改变任一固定温度 \(q\) 上的两点极限残差、任一 \(\beta\)-Rényi 极限、任一 Csisz\'ar 极限、或 escort Fisher 曲率。
\end{corollary}

\begin{proof}
该变形既保持某个固定 \(\alpha>1\) 的极限幂散度不变，由定理 \ref{thm:conclusion-binfold-single-powerdiv-asymptotic-statistical-completeness}，便必须保持 \(\varphi\) 不变。而上述全部极限对象都已知是 \(\varphi\) 的显式函数，故它们逐一不变。于是该变形在极限实验层面不再携带独立参数。证毕。
\end{proof}


\noindent
因此，这条二态链在结论层形成如下闭合结构：末位统计量首先以显式核实现整个 escort 族的渐近 Blackwell 等价；相应极限模型在平方根坐标下是一段可完全参数化、且总长度恰为 \(2\arctan(\varphi^{-1/2})\) 的有限 Fisher--Rao 圆弧；全部 escort--escort 的 Csiszár 几何及其总变差/\(\chi^2\) 截面都由同一两态骨架承载，而相对均匀基线的 \(\chi^2\)-型黄金识别则由同一骨架的另一截面给出；更进一步，任意一个固定 \(\alpha>1\) 上的幂散度极限标量已经足以反演 \(\varphi\)，并回收整条 escort 温度族的 KL 极限与两尺度残差律，甚至把所有保持该标量不变的变形压缩为同一个 Le Cam 极限实验；在最终的决策论层面，任意有界损失下的最优 Bayes 风险也全部塌缩为同一 Bernoulli 实验的风险函数。由此，bin-fold escort 族的极限结构不再只是若干分散的闭式公式，而成为一个同时封闭于信息几何、散度理论与统计决策的单比特实验模型。

\endinput
